\section{Entrenamiento de RNN: retropropagación y problemas de gradiente}

\subsection{Función de pérdida de la RNN}

Al igual que en las redes feedforward, entrenar una RNN también requiere definir una función de pérdida. La diferencia es que la RNN puede calcular un valor de pérdida en \textbf{cada paso de tiempo}:

$$\mathcal{L}_t = \text{Loss}(\hat{y}_t, y_t)$$

La pérdida total es la suma de las pérdidas de todos los pasos de tiempo:

$$\mathcal{L} = \sum_{t=1}^{T} \mathcal{L}_t$$

\subsection{Backpropagation Through Time (BPTT)}

Entrenar una RNN utiliza una versión extendida de la retropropagación: la \textbf{retropropagación a través del tiempo} (Backpropagation Through Time, BPTT). Su idea central es:

\begin{enumerate}
    \item Propagación hacia adelante: calcular el estado oculto y la salida paso a paso en el tiempo
    \item Calcular la pérdida en cada paso de tiempo
    \item Retropropagación: propagar el error a lo largo del eje temporal de atrás hacia adelante, calculando el gradiente de cada parámetro
    \item Actualizar los parámetros
\end{enumerate}

\begin{figure}[H]
\centering
\includegraphics[width=0.85\textwidth]{slides-images/slide-45.jpg}
\caption{Flujo del gradiente en una RNN estándar\protect\footnotemark}
\end{figure}
\footnotetext{Diapositiva 45.}

\subsection{Desvanecimiento y explosión del gradiente}

Durante el proceso de BPTT, el gradiente debe pasar por múltiples multiplicaciones de matrices (multiplicaciones sucesivas de $W_{hh}$) y por el producto sucesivo de las derivadas de la función de activación. Esto provoca dos problemas numéricos graves:

\begin{importantbox}{Desvanecimiento y explosión del gradiente}
\begin{itemize}
    \item \textbf{Explosión del gradiente} (Exploding Gradients): cuando el valor propio de $W_{hh}$ es $> 1$, el gradiente crece exponencialmente con los pasos de tiempo, lo que produce actualizaciones de parámetros excesivas e inestabilidad en el entrenamiento
    \item \textbf{Desvanecimiento del gradiente} (Vanishing Gradients): cuando el valor propio de $W_{hh}$ es $< 1$, el gradiente decae exponencialmente con los pasos de tiempo hasta acercarse a cero, lo que impide que la red aprenda dependencias de largo plazo
\end{itemize}
\end{importantbox}

\begin{figure}[H]
\centering
\includegraphics[width=0.85\textwidth]{slides-images/slide-48.jpg}
\caption{El problema del desvanecimiento del gradiente: calcular el gradiente respecto de $h_0$ implica multiplicaciones sucesivas de $W_{hh}$\protect\footnotemark}
\end{figure}
\footnotetext{Diapositiva 48.}

\subsection{La dificultad de las dependencias de largo plazo}

El problema del desvanecimiento del gradiente hace que a la RNN le resulte difícil capturar dependencias de largo plazo:

\begin{figure}[H]
\centering
\includegraphics[width=0.85\textwidth]{slides-images/slide-52.jpg}
\caption{El desvanecimiento del gradiente hace que la RNN tienda a capturar dependencias de corto plazo\protect\footnotemark}
\end{figure}
\footnotetext{Diapositiva 52.}

\begin{warningbox}{Desvanecimiento del gradiente $\neq$ gradiente igual a cero}
El desvanecimiento del gradiente no significa que el gradiente desaparezca por completo, sino que la señal de gradiente proveniente de \textbf{pasos de tiempo más lejanos} se vuelve extremadamente débil, de modo que la actualización de los parámetros del modelo queda dominada casi por completo por el gradiente de los pasos de tiempo recientes. Por ello, el modelo tiende a aprender dependencias de corto plazo y le cuesta capturar patrones de largo plazo.
\end{warningbox}

\subsection{Estrategias para resolver los problemas de gradiente}

Existen tres grandes tipos de métodos para resolver el desvanecimiento/explosión del gradiente:

\begin{enumerate}
    \item \textbf{Elección de la función de activación}: usar funciones de activación como ReLU, que no comprimen el gradiente en la zona de saturación
    \item \textbf{Inicialización de los pesos}: diseñar cuidadosamente la estrategia de inicialización de los pesos para evitar que el gradiente se escale rápidamente durante la propagación
    \item \textbf{Mejoras en la arquitectura de red}: diseñar mecanismos de compuertas específicos para controlar el flujo de información, como LSTM y GRU
\end{enumerate}

\subsection{LSTM: redes de memoria de largo y corto plazo}

\textbf{Long Short-Term Memory (LSTM)} es la variante de RNN más conocida para resolver el problema del desvanecimiento del gradiente. Su idea central es agregar \textbf{mecanismos de compuertas} (Gates) dentro de la unidad estándar de la RNN, con los que se controla de manera selectiva el olvido, el almacenamiento y la salida de información.

\begin{importantbox}{Las tres compuertas de la LSTM}
\begin{itemize}
    \item \textbf{Compuerta de olvido} (Forget Gate): decide qué información antigua debe descartarse
    \item \textbf{Compuerta de entrada} (Input Gate): decide qué información nueva debe almacenarse
    \item \textbf{Compuerta de salida} (Output Gate): decide qué información debe emitirse hacia el siguiente paso de tiempo
\end{itemize}
El valor de estas compuertas se controla mediante la función sigmoide, cuya salida está entre $[0, 1]$ y representa desde «olvido completo» hasta «retención completa».
\end{importantbox}

Mediante este mecanismo de compuertas, la LSTM permite que el gradiente se propague entre pasos de tiempo a lo largo de una especie de «autopista» (Cell State), lo que alivia en gran medida el problema del desvanecimiento del gradiente y permite que la red aprenda dependencias de más largo plazo.

\begin{knowledgebox}{Importancia histórica de la LSTM}
La LSTM fue propuesta por Hochreiter y Schmidhuber en 1997, y es una de las innovaciones arquitectónicas más importantes en la historia del aprendizaje profundo. Antes de la aparición del Transformer (en 2017), la LSTM y sus variantes (como la GRU) fueron la opción dominante para procesar datos secuenciales, con un uso extendido en la traducción automática, el reconocimiento de voz, la generación de texto y otras áreas.
\end{knowledgebox}

\subsection{Resumen de la sección}

El reto central al entrenar una RNN es el problema del desvanecimiento/explosión del gradiente, que limita la capacidad de la RNN estándar para aprender dependencias de largo plazo. La LSTM alivia eficazmente este problema mediante su mecanismo de compuertas, pero en esencia sigue procesando la secuencia paso a paso en el tiempo, con limitaciones tanto en la eficiencia computacional como en la capacidad de modelar secuencias largas.

\newpage
